\chapter{Protocolo experimental: de la señal a la cartera}
\label{chap:diseno-experimental}

El capítulo anterior describió cómo se produce una ordenación. Éste explica cómo se decide qué
configuración usar y cómo esa ordenación se convierte en posiciones, sin que en ningún momento la
rentabilidad observada influya en las decisiones predictivas.

El orden de exposición reproduce el orden real del procedimiento: primero se elige la configuración
predictiva comparando Rank-IC, después se congela, y sólo entonces se traduce el ranking a alfa
esperado y a una cartera. La última parte del capítulo explica cómo se optimizó esa cartera y por
qué ese segundo procedimiento usa una búsqueda distinta de la primera.

Conviene precisar qué clase de cifras aparecen en este capítulo y cuáles no, porque la distinción
organiza el resto del documento. Aquí se presentan las \emph{trazas de decisión}: qué alternativa se
comparó con cuál, cuánta ventaja pareada midió cada una y qué regla resolvió el empate. Son
necesarias, porque sin ellas la afirmación de que el protocolo se aplicó tal como se declara no
sería comprobable, y un protocolo no auditable no es un protocolo.

Lo que no aparece aquí es la interpretación de esas cifras como evidencia sobre el mundo. Que una
alternativa midiera peor documenta una decisión; si la señal resultante es real, si resiste
contrastes adversariales y si se traduce en rentabilidad son preguntas distintas, y se responden con
otros datos en los Capítulos~\ref{chap:resultados} y~\ref{chap:resultados-economicos}.

\section{Por qué el protocolo se aplicó tres veces}
\label{sec:cadena-studies}

El protocolo que describe este capítulo no se ejecutó una vez sino tres, encadenadas, y conviene
explicarlo antes de entrar en la mecánica, porque todas las trazas que siguen son las de la tercera
pasada --- a la que el documento llama \emph{pasada de referencia}.

La selección secuencial recorre las fases en orden y cada una hereda al incumbente de la anterior.
Eso significa que el resultado depende del punto de partida: una variable decidida en la primera
fase se decide con los valores por defecto de todas las demás, y no vuelve a revisarse aunque las
decisiones posteriores cambien el contexto en que aquella elección tenía sentido. La primera pasada
lo dejó claro en cifras: su ganador se apartaba en \textbf{ocho variables} del baseline recomendado
del catálogo, de modo que la mayoría de sus decisiones se habían tomado partiendo de una
configuración que el propio estudio acabó descartando.

La respuesta fue encadenar los estudios. El ganador de cada pasada se convierte en el baseline de la
siguiente, que vuelve a recorrer las fases desde ese punto. Es un ascenso por coordenadas: cada
pasada sólo puede mejorar o confirmar, porque el incumbente se conserva cuando no hay evidencia
suficiente para desplazarlo. La primera modificó ocho variables respecto de su punto de partida; la
segunda, una sola (\artefacto{meta\_method}, de acotado a libre); la tercera, dos
(\artefacto{execution\_lag\_days} y \artefacto{target\_size}). El procedimiento se queda sin cambios
que hacer, que es lo que debe ocurrir si converge en lugar de explorar al azar, y es la razón de
detenerse en tres. Las Tablas~\ref{tab:cadena} y~\ref{tab:cadena-config} del
Capítulo~\ref{chap:resultados} recogen qué obtuvo y qué cambió cada pasada.

Conviene declarar el coste metodológico, porque no es gratuito. Encadenar estudios multiplica el
número de configuraciones evaluadas sobre la misma ventana de selección, y por tanto aumenta la
probabilidad de que la mejor cifra observada deba parte de su tamaño a la propia búsqueda. La cadena
no es una forma de evitar el sobreajuste: es una forma de sistematizarlo y hacerlo visible. Lo que
protege al resultado no es el procedimiento de búsqueda sino la ventana reservada, que no participa
en ninguna de las tres pasadas.

\section{Optimización secuencial y regla de decisión}

El catálogo cerrado se recorre en cuatro fases: temporal, representación, modelo y meta. Se parte
de un \textit{baseline}; en cada fase se modifica una variable y se comparan sus candidatos con el
incumbente acumulado. Así, el coste científico es una suma de alternativas, no una búsqueda
cartesiana opaca. Las 71 configuraciones ejecutadas en la pasada de referencia quedan registradas en
\artefacto{decisions.json} junto con cada elección, y las de las dos pasadas anteriores conservan
sus propios registros.

Un candidato debe tener observaciones suficientes, no caer por debajo de $-0,02$ en ninguna era
disponible y superar una puerta pareada: dominar al incumbente (ventaja media positiva y mejor en
más de la mitad de cohortes) o ser no inferior con límite inferior bootstrap al 90\% superior a
$-0,01$. Las diferencias se calculan por periodo mensual, lo que permite comparar lags distintos
sin confundir el desplazamiento de sus fechas de ejecución con ausencia de evidencia.

\begin{figure}[H]
  \centering
  \begin{tikzpicture}[node distance=5mm, >=Latex, every node/.style={font=\small,align=center}]
    \node[draw,rounded corners,fill=blue!8,minimum width=27mm,minimum height=10mm] (t) {Temporal};
    \node[draw,rounded corners,fill=teal!10,right=of t,minimum width=30mm,minimum height=10mm] (r) {Representación};
    \node[draw,rounded corners,fill=orange!12,right=of r,minimum width=25mm,minimum height=10mm] (m) {Modelo};
    \node[draw,rounded corners,fill=green!10,right=of m,minimum width=25mm,minimum height=10mm] (x) {Meta};
    \node[draw,rounded corners,fill=gray!15,below=8mm of r,minimum width=48mm,minimum height=10mm] (w) {Ganador congelado};
    \node[draw,rounded corners,fill=yellow!18,right=of w,minimum width=38mm,minimum height=10mm] (d) {Robustez y cartera\\diagnósticas};
    \draw[->,thick] (t)--(r); \draw[->,thick] (r)--(m); \draw[->,thick] (m)--(x);
    \draw[->,thick] (x.south)--++(0,-5mm)-|(w.north);
    \draw[->,thick] (w)--(d);
  \end{tikzpicture}
  \caption{Protocolo de selección. Sólo las cuatro fases superiores pueden modificar el ganador;
  la evidencia económica se calcula una vez congelada la señal.}
  \label{fig:protocolo}
\end{figure}

La Tabla~\ref{tab:decisiones} hace auditable esa traza para la pasada de referencia. Conviene
explicar qué columna se tabula y por qué, porque la elección no es inocente. La representación
habitual --- el Rank-IC acumulado tras cada paso, en forma de escalera --- sería aquí una columna
casi constante, por los motivos que desarrolla el apartado siguiente. Lo informativo, entonces, no
es el valor que se conserva sino \textbf{la alternativa que se descarta}: cuánto Rank-IC habría
costado adoptarla. Esa es la columna que se tabula, ordenada de mayor a menor coste.

\begin{table}[H]
  \centering
  \small
  \caption{Qué alternativa se rechazó en cada decisión y cuánto habría costado adoptarla.}
  \label{tab:decisiones}
  \begin{tabular}{lllrrl}
\toprule
Variable & Ganador & Mejor alternativa & Su Rank-IC & Coste & Regla \\
\midrule
neutralize\_by\_sector & False & True & 0,0576 & -0,0429 & Rank-IC robusto \\
recency\_weighting & off & linear & 0,0632 & -0,0373 & Rank-IC robusto \\
target\_horizon\_months & 12 & 6 & 0,0684 & -0,0322 & Rank-IC robusto \\
feature\_preset & all & core & 0,0716 & -0,0290 & Rank-IC robusto \\
snapshot\_step\_months & 1 & 3 & 0,0916 & -0,0089 & Rank-IC robusto \\
meta\_method & stacked rolling free & stacked rolling bounded & 0,1005 & -0,0085 & Rank-IC robusto \\
train\_lookback\_years & 8 & 12 & 0,0953 & -0,0052 & Rank-IC robusto \\
feature\_weighting\_mode & oos stability prune & model native & 0,0959 & -0,0047 & Rank-IC robusto \\
fundamental\_momentum & True & False & 0,0967 & -0,0038 & Rank-IC robusto \\
winsorization & 0.0 & 0.01 & 0,0984 & -0,0021 & Rank-IC robusto \\
lgbm\_max\_depth & 3 & 4 & 0,0992 & -0,0014 & Rank-IC robusto \\
lgbm\_learning\_rate & 0.03 & 0.05 & 0,0995 & -0,0010 & Rank-IC robusto \\
lgbm\_n\_estimators & 100 & 200 & 0,0997 & -0,0008 & Rank-IC robusto \\
execution\_lag\_days & 60 & 30 & 0,1000 & -0,0005 & Empate: simplicidad \\
max\_features\_per\_agent & 20 & 12 & 0,1000 & -0,0005 & Rank-IC robusto \\
market\_regime\_feature & False & True & 0,1001 & -0,0004 & Rank-IC robusto \\
lgbm\_min\_child\_samples & 20 & 50 & 0,1009 & 0,0004 & Empate: simplicidad \\
\bottomrule
\end{tabular}

  \caption*{Fuente: \artefacto{decisions.json}. El coste es el Rank-IC de la mejor alternativa menos
  el del ganador; un coste \emph{positivo} indica que la alternativa medía mejor y que el ganador se
  decidió por otra vía, lo que sólo ocurre en los empates técnicos. No debe confundirse con la
  \emph{ventaja pareada}, que promedia la diferencia cohorte a cohorte y por eso no coincide
  exactamente: en \artefacto{meta\_method}, por ejemplo, el coste es $-0{,}0085$ y la ventaja
  pareada $+0{,}00827$.}
\end{table}

Leída así, la tabla separa con claridad tres grupos. Cuatro decisiones eran caras de equivocar:
neutralizar por sector habría costado $-0{,}0429$ de Rank-IC, aplicar pesos de recencia $-0{,}0373$,
acortar el horizonte a seis meses $-0{,}0322$ y reducir el conjunto de variables a \texttt{core}
$-0{,}0290$. Son las decisiones estructurales, y las cuatro se resolvieron con holgura.

Un segundo grupo, el más numeroso, tiene costes por debajo de 0,005: los hiperparámetros de LightGBM
(profundidad, número de árboles, tasa de aprendizaje), la winsorización o el número máximo de
variables por agente. Ahí la evidencia apenas distingue entre alternativas, y conviene decirlo
abiertamente: el sistema no debe su capacidad predictiva al ajuste fino del modelo.

El tercer grupo son las dos decisiones resueltas por empate técnico, que la última columna marca de
forma explícita. En \artefacto{execution\_lag\_days} la alternativa de 30 días quedó $0{,}0005$ por
debajo del ganador y en \artefacto{lgbm\_min\_child\_samples} el valor 50 medía $0{,}0004$
\emph{mejor}. Es decir, en el segundo caso el ganador no es el valor que mejor midió. Ninguna de las
dos debe presentarse como hallazgo empírico, y por eso figuran con su regla al lado.

\subsection*{Por qué la escalera es plana}

El resultado de la tercera pasada es contundente y va en contra de la intuición habitual sobre este
tipo de búsquedas: \textbf{las diecisiete decisiones confirmaron su valor de partida}, dieciséis de
ellas con una ventaja pareada exactamente nula. La única comparación que midió una diferencia
apreciable fue \artefacto{meta\_method}, donde la variante libre obtuvo 0,1090 frente a 0,1005 de la
acotada; también ahí el procedimiento conservó el valor que ya traía.

Esa planitud es exactamente lo que cabía esperar de una cadena que ha convergido: la pasada de
referencia parte del ganador de la anterior, de modo que casi todas las variables llegan ya en su
mejor valor conocido y la comparación pareada no encuentra motivo para moverlas. La primera pasada,
que partía del baseline del catálogo, presentaba un perfil muy distinto.

Este patrón es una defensa parcial frente a la objeción de sobreajuste por multiplicidad. Un
protocolo que aceptase cualquier mejora aparente habría desplazado al incumbente muchas veces,
persiguiendo diferencias dentro del ruido. Que las diecisiete decisiones terminen en «no cambiar»
indica que las puertas de elegibilidad están operando. No elimina el problema --- el
Deflated Sharpe del Capítulo~\ref{chap:resultados} sigue penalizando por las configuraciones
ensayadas, y la cadena eleva ese recuento en lugar de reducirlo ---, pero acota su alcance: la
multiplicidad afecta sobre todo a la afirmación económica, no a la traza de decisiones predictivas.

La regla de simplicidad que resolvió esas dos decisiones opera bajo un umbral de empate técnico de
0,002 en ventaja pareada. Es una preferencia declarada de antemano --- ante evidencia equivalente,
se elige la opción más simple o la más conservadora --- y no una medición, lo que delimita
exactamente qué puede afirmarse sobre los valores que produce.

\subsection*{El catálogo cerrado como garantía}

La afirmación de que nada se optimizó fuera del catálogo sólo es verificable si el catálogo es
público. El Anexo~\ref{anexo:catalogo} reproduce las treinta y tres variables con su rejilla
completa, su etapa, si son predictivas y el valor que tomó el ganador. Su función en este documento
no es descriptiva sino probatoria: cualquier lector puede comprobar que el valor ganador de cada
variable pertenece a su rejilla declarada, y que ninguna variable ajena aparece en la configuración
final.

Doce de esas treinta y tres variables son de cartera y están marcadas como no predictivas. La
distinción es la que sostiene el Capítulo~\ref{chap:resultados}: se ejecutan después de congelar el
ganador, se reportan como diagnóstico y no pueden modificarlo. Es lo que convierte el barrido de la
Sección~\ref{sec:cartera-rejilla} en un experimento controlado en lugar de una segunda ronda de
optimización.

\section{Por qué la selección predictiva es secuencial}
\label{sec:seleccion-secuencial}

Una búsqueda cartesiana de todos los valores posibles permitiría encontrar combinaciones muy
ajustadas a la historia, dificultaría explicar de dónde surge el ganador y multiplicaría las
pruebas. El protocolo fija un orden: temporalidad, representación, modelo y meta. Cada fase hereda
el incumbent anterior, de modo que una decisión posterior no puede reabrir silenciosamente una
elección temporal. Si dos alternativas difieren menos de 0,002 en ventaja pareada, se prefiere la
más simple; si tampoco hay evidencia suficiente, se conserva el incumbent.

Conviene precisar el alcance de este argumento, porque el trabajo emplea después una búsqueda
cartesiana y sería contradictorio no explicar por qué. El razonamiento anterior vale para el plano
\textbf{predictivo}, y vale por dos motivos que no son intercambiables. El primero es estadístico:
la selección predictiva se decide sobre el Rank-IC, la misma magnitud que después se reporta como
resultado principal, de modo que cada comparación adicional infla el mejor valor observado y
contamina precisamente la cifra que sostiene el trabajo. El segundo es expositivo: el objetivo del
Model Study no es solo encontrar una configuración buena, sino poder explicar de dónde sale cada
decisión, y una rejilla de miles de puntos devuelve un ganador que nadie puede narrar.

Ninguno de los dos motivos se aplica al plano de la \textbf{cartera}. Las reglas de cartera no
tocan el modelo ni modifican el Rank-IC: operan sobre una señal ya congelada, así que optimizarlas
no puede inflar la métrica predictiva que el documento reporta. Y su elección no necesita narrarse
como una escalera de decisiones porque no lo es: las variables de cartera interactúan entre sí, y un
protocolo secuencial no solo sería innecesario sino directamente inadecuado, como se argumenta en el
Anexo~\ref{anexo:desarrollo}. Queda un riesgo real, la multiplicidad de comparaciones, que
no se disimula: se declara en el Capítulo~\ref{chap:limitaciones} y se acota por construcción
impidiendo que la rejilla llegue a calcular la era reservada.

\section{Ejemplo de comparación pareada}

\begin{ejemplo}[Mecánica de la comparación]
Un candidato y el incumbente producen Rank-IC sobre las mismas 40 cohortes. En cada fecha se resta
el Rank-IC del incumbente al del candidato y se re-muestrean bloques de doce cohortes para preservar
la dependencia inducida por el horizonte anual. Si la ventaja media resultase 0,0208, con el
candidato mejor en el 57,5\,\% de las cohortes y un intervalo \textit{bootstrap} al 90\,\% de
$[0{,}0108;\,0{,}0372]$, habría evidencia de dominancia.
\end{ejemplo}

Las cifras del ejemplo son ilustrativas, pero corresponden al orden de magnitud de un caso real: así
se decidió la cadencia mensual frente al baseline trimestral en la primera pasada de la cadena, con
una ventaja pareada de $+0{,}0208$.

La rentabilidad de cartera no influye en este paso. Un CAGR alto puede surgir de concentración,
exposición accidental a una acción extrema o efectivo, sin demostrar que se ordene mejor el
universo. El horizonte de seis meses ilustra un rechazo medido: su ventaja pareada fue $-0,0265$ y
el intervalo $[-0,0552;-0,0109]$. Mostrar este resultado evita presentar la selección como una
confirmación automática de la configuración preferida.

\section{De rango a alfa esperado}

Todo lo anterior ocurre en el plano predictivo y termina con una configuración congelada. A partir
de aquí el problema cambia de naturaleza: ya no se trata de ordenar mejor, sino de convertir una
ordenación en posiciones concretas sin destruir por el camino la ventaja que contiene.

El primer paso es que la cartera necesita una magnitud económica, no sólo un percentil. Sobre cohortes cerradas se estima
una relación entre ventiles de \texttt{meta\_rank} y retorno excedente anualizado. Se usan veinte
grupos para no ajustar una curva a medias de muy pocas acciones, aunque la función se evalúa sobre
el rango continuo. Si la pendiente no es creciente, la evidencia se amplía en cascada a dieciséis
trimestres y al histórico completo; la salvaguarda final queda registrada como supuesto y no se
presenta como estimación empírica.

\section{Construcción de cartera}

Seis variables gobiernan cómo se traduce el ranking en posiciones: cuántas se mantienen
(\artefacto{target\_size}), cómo se reparte el peso entre ellas (\artefacto{sizing\_mode}), cuánto
efectivo se tolera (\artefacto{max\_cash\_weight}), cuánto debe retenerse una posición
(\artefacto{minimum\_holding\_period}), por debajo de qué percentil se vende
(\artefacto{coverage\_percentile\_floor}) y con qué frecuencia se corrige la deriva de los pesos
(\artefacto{rebalance\_drift\_tolerance}). El catálogo declara la rejilla de cada una; qué valores
toman finalmente se decide con el procedimiento de la Sección~\ref{sec:portfolio-study} y se
presenta en el Capítulo~\ref{chap:resultados-economicos}.

Las restantes condiciones de simulación sí son constantes del entorno y no se optimizan: comisión de
5 pb, \textit{slippage} de 10 pb y SPY como benchmark, nunca como posición. Los umbrales se expresan
en puntos básicos anuales de alfa esperado para poder compararlos entre horizontes. Una rotación
requiere que la diferencia esperada cubra el coste de ida y vuelta y un margen de 50 pb; una compra
nueva usa histéresis. Estas reglas pretenden preservar señal y controlar rotación, no volver a
optimizar el modelo por rentabilidad.

La era 2025--2026 fue aislada antes de aplicar la selección. Sus observaciones no participan ni en
los pesos del meta ni en la elección de parámetros; su papel es exclusivamente de confirmación
posterior.

\section{Ejemplo de decisión y coste}

\begin{ejemplo}[Cuándo se rota y cuándo no]
Con horizonte de doce meses, comisión de 5 pb y \textit{slippage} de 10 pb, una rotación completa
cuesta 30 pb. Si la peor posición tuviese alfa esperado 60 pb, una candidata externa 250 pb y el
margen de rotación fuese 50 pb, la mejora de 190 pb superaría los 80 pb exigidos y se rotaría. Si la
candidata ofreciese 120 pb, la mejora no cubriría coste y margen, y se conservaría la posición.
\end{ejemplo}

El ranking ordena preferencias; la operación necesita además una ventaja neta. Las cifras anteriores
son inventadas para ilustrar la regla: los costes sí son los de la simulación, pero los alfas no
proceden de ninguna medición.

Las entradas nuevas incorporan histéresis y el efectivo rinde 0\%. Bajo la política
\texttt{opportunity\_cash}, una plaza puede permanecer en efectivo sólo si respeta el tope declarado
y el suelo de diversificación; con el tope fijado en cero, esa misma plaza debe reinvertirse
necesariamente. Un mínimo de tenencia, si se activa, bloquea ventas por rotación o caída de alfa
hasta cumplir la fracción de horizonte configurada. Estas reglas controlan fricción y estabilidad;
no pueden alterar el ganador predictivo.

\section{El Portfolio Study: optimizar la cartera sin tocar la señal}
\label{sec:portfolio-study}

Las seis variables anteriores tienen que tomar algún valor, y elegirlo por su rentabilidad conocida
parecería contradecir la regla que gobierna todo el protocolo. No lo hace, y conviene ver por qué,
porque de ese razonamiento sale el diseño del cuarto estudio.

Lo que la regla protege es la \emph{métrica predictiva}: si se elige una configuración porque su
Rank-IC es alto y después se reporta ese mismo Rank-IC como resultado, la cifra está inflada por la
propia búsqueda. Las reglas de cartera no participan de ese circuito. Operan sobre una señal ya
congelada, no reentrenan nada y no pueden modificar el Rank-IC. Lo que sí hay que protegerles es
otra cosa: la era reservada. Y eso no se protege renunciando a optimizar, sino con una
construcción.

De ahí el \emph{Portfolio Study}. Toma el ganador del último Model Study, no reentrena nada, y
recorre el producto cartesiano de las seis variables que gobiernan la cartera midiendo cada
combinación sobre una serie de puntuaciones \textbf{recortada en 2024}. La distinción es fina pero
decisiva: no se trata de calcular el resultado de cada cartera en la era reservada y después
ignorarlo --- si existiera, bastaría con mirarlo --- sino de que ese resultado \emph{no llega a
calcularse}. Sólo la ganadora, ya elegida, se reevalúa después sobre la serie completa, y esa
reevaluación es una comprobación, nunca un criterio.

\subsection*{Qué es el Information Ratio y por qué se optimiza éste}

El criterio de selección deja de ser el exceso geométrico y pasa a ser el \emph{Information Ratio}:
el exceso medio sobre el índice dividido por la volatilidad de ese exceso. Premia la
\textbf{consistencia}, no el alfa bruto. Una cartera que bate al índice un 2\,\% todos los años
tiene un IR altísimo; otra que alterna $+15$\,\% y $-11$\,\% puede tener el mismo exceso medio y un
IR pésimo. Para un trabajo cuya pregunta es si una señal se sostiene fuera de muestra, la segunda
cartera no es equivalente a la primera aunque el promedio coincida, y optimizar por exceso bruto no
sabría distinguirlas.

Las seis variables mueven ese cociente por vías distintas, y conviene enunciarlas porque explican la
forma de los resultados:

\begin{itemize}
  \item \artefacto{target\_size} gobierna el compromiso central. Pocas posiciones concentran la
  señal pero hacen el exceso volátil; muchas lo estabilizan a costa de parecerse al índice, y una
  cartera que se parece al índice tiene poco exceso que mostrar en el numerador.
  \item \artefacto{max\_cash\_weight} amortigua numerador y denominador a la vez. El efectivo no se
  remunera en esta simulación, de modo que sostener caja resta rentabilidad esperada, pero también
  resta varianza del exceso.
  \item \artefacto{sizing\_mode} decide si el peso sigue al alfa esperado o se reparte por igual.
  Concentrar en las convicciones altas sube el exceso y la varianza simultáneamente.
  \item \artefacto{minimum\_holding\_period} reduce fricciones, pero no de forma monótona: retener
  de más impide corregir cuando la señal cambia de opinión, y ese coste de oportunidad acaba
  superando al ahorro en costes.
  \item \artefacto{coverage\_percentile\_floor} corta la cola inferior del ranking, forzando la
  venta de lo que cae por debajo del percentil declarado.
  \item \artefacto{rebalance\_drift\_tolerance} es pura fricción: no cambia qué se tiene, sino cada
  cuánto se corrige el peso de lo que ya se tiene.
\end{itemize}

\subsection*{Por qué cartesiano y no secuencial}

La Sección~\ref{sec:seleccion-secuencial} argumentó que la selección predictiva debe ser secuencial.
Aquí se hace lo contrario, y no por comodidad: las seis variables \textbf{interactúan},
y un procedimiento que las recorriera de una en una obtendría una respuesta que depende del orden en
que se pregunte.

Dos ejemplos concretos lo muestran. El suelo efectivo de diversificación no lo fija ninguna variable
por separado, sino \artefacto{target\_size} y \artefacto{max\_cash\_weight} conjuntamente: el peso
máximo por posición depende de cuántas posiciones caben y de cuánta caja se tolera. Y
\artefacto{coverage\_percentile\_floor} hace cosas \emph{opuestas} según el tope de efectivo: con
tope cero, la plaza que deja una posición vendida por suelo se recompra necesariamente, de modo que
el suelo actúa como un mecanismo de sustitución; con tope positivo, esa misma plaza puede quedarse
en caja, y entonces el suelo actúa como un mecanismo de reducción de exposición. La misma variable,
con el mismo valor, significa dos cosas distintas según el valor de otra.

Un ascenso por coordenadas sobre variables que interactúan así converge a un óptimo local que
depende del punto de partida. El cartesiano no tiene ese problema. Tiene otro, la multiplicidad de
comparaciones, que se declara abiertamente en el Capítulo~\ref{chap:limitaciones}.

\subsection*{Los perfiles quedan fuera de la rejilla}

Los ocho perfiles de inversor no compiten en el cartesiano, y la razón es que no son variables de
cartera. Un perfil \emph{reordena la señal}: sustituye el ranking del meta-agente y recalibra el
alfa esperado. Opera, por tanto, en el plano predictivo y no en el de gestión. Meterlos en la
rejilla equivaldría a elegir el estilo de inversor por su rentabilidad conocida, que es justamente
lo que la regla protege.

Lo que sí se hace es reevaluar los ocho con la cartera ya ganadora, para responder a una pregunta
legítimamente distinta: cómo le habría ido a cada estilo de inversor con la mejor gestión
disponible. Es un diagnóstico informativo y se presenta como tal; \textbf{ningún perfil se
selecciona por su Information Ratio}.

\section*{Qué queda fijado al terminar este capítulo}

Al concluir el protocolo existen dos objetos congelados y ninguna cifra de resultado interpretada.
El primero es una \emph{configuración predictiva}: un valor para cada variable del catálogo,
elegido comparando Rank-IC pareado y sin que la rentabilidad interviniera en ninguna comparación.
El segundo es un \emph{procedimiento de cartera}: seis variables con su rejilla, un criterio de
optimización --- el Information Ratio --- y la garantía constructiva de que la era reservada no
llega a calcularse durante la búsqueda.

Con eso, las dos preguntas empíricas del trabajo quedan planteadas en el orden en que pueden
responderse. La primera es si la ordenación que ese sistema produce existe realmente y resiste los
intentos de explicarla por azar; es independiente de cualquier cartera y la responde el
Capítulo~\ref{chap:resultados}. La segunda sólo tiene sentido si la primera se responde
afirmativamente: si esa ordenación sobrevive al convertirse en posiciones. La responde el
Capítulo~\ref{chap:resultados-economicos}.
